\section{Gestion de projet}

$ \ast $ \textit{La section suivante pourra être négociée de manière à ce que toutes les parties prenantes soient en accord avec le contenu.}

\subsection{Décomposition en work packages \gloss{wp}}
Le projet est décomposé en différents work packages \gloss{wp}, ce qui permet une répartition claire des tâches. 
Chaque collaborateur travaille sur une partie distincte du projet, assurant ainsi une organisation efficace et une meilleure gestion des responsabilités.

\subsubsection{Fonctionnalités principales}
Cette section présente les fonctionnalités principales du projet, qui sont essentielles pour atteindre les objectifs fixés et garantir le succès du projet. \\
\begin{itemize}
    \item \textbf{Hardware} : Jules Noumba s'occupera de l'agencement sur le piège caméra et des différents capteurs qui permettront de récupérer les \gloss{metadonnees} liées aux photos.
    \item \textbf{Énergie} : Nolan Buchet aura pour rôle de gérer les consommations d'énergie des composants agencés par Jules Noumba.
    \item \textbf{Transmission} : Joris Menard s'occupera de la transmission sans fil des données entre le piège caméra et l'unité de stockage distante.
    \item \textbf{Stockage des données} : Davy Clark s'occupera de trouver une solution pour stocker les données récupérées par le piège caméra et de la développer.
    \item \textbf{Module caméra} : Xueying Xu aura pour rôle de choisir et de mettre en place le module caméra afin de capturer des images de bonne qualité, facilitant ainsi l'identification si cette fonctionnalité est développée.
    \item \textbf{Gestion de projet} : En tant que chef de projet, Ewen Le Callet sera chargé de gérer l'équipe, de répartir les tâches, de vérifier la qualité du travail réalisé et de s'assurer que les livrables sont rendus dans les délais.
    \item \textbf{Communication avec le client} : Ewen Le Callet aura également le rôle d'interlocuteur entre l'équipe de projet et le client afin de clarifier certains points ou de tenir compte de l'avancée du projet.
\end{itemize}

\subsubsection{Fonctionnalités secondaires}
Cette section décrit les fonctionnalités secondaires du projet, qui complètent les fonctionnalités principales et ajoutent de la valeur à l'ensemble du projet. \\
\begin{itemize}
    \item \textbf{interface homme machine (ihm)} : Xueying Xu aura le rôle de développer une interface homme machine \gloss{ihm}, permettant à l'utilisateur du produit final de visualiser les données récupérées par le piège caméra.
    \item \textbf{IA} : Ewen Le Callet s'occupera de mettre en place l'intelligence artificielle afin de réaliser une identification des animaux sur les photos prises par le piège caméra.
\end{itemize}

\subsection{Signatures du projet}
\begin{table}[h!]
\centering
\begin{tabularx}{\textwidth}{|X|X|X|}
\hline
\textbf{Date et signature du client} &
\textbf{Date et signature du chef de projet} &
\textbf{Date et signature des coachs de projet} \\
\hline
\makebox[0pt][l]{\rule{0pt}{2cm}} & 
\makebox[0pt][l]{\rule{0pt}{2cm}} & 
\makebox[0pt][l]{\rule{0pt}{2cm}} \\
\hline
\end{tabularx}
\end{table}
